Inverse folding---finding sequences that adopt a prescribed structure---is now dominated by conditional autoregressive networks such as ProteinMPNN~\cite{dauparas2022robust} and ESM-IF1~\cite{hsu2022learning}, whose designs are usually judged by refolding them with a structure predictor such as AlphaFold~\cite{jumper2021highly}. That evaluation is approximate (the predictor can be wrong), expensive, and conflates two questions: whether the generator proposes good sequences, and how many verified candidates it needs. The classical alternative is stochastic search against an energy-like objective, for example simulated annealing~\cite{kirkpatrick1983optimization}. Which of the two is preferable depends on the cost of the oracle that decides whether a candidate really folds to the target.

The two-dimensional hydrophobic--polar (HP) lattice model~\cite{lau1989lattice} offers a setting where this question can be asked with an \emph{exact} oracle. Chains of $N=16$ residues have only \rhval{data/enumeration/hp16/n_conformations/mean} distinct conformations modulo rotation and reflection, so the ground state and its degeneracy can be computed for every one of the $2^{16}$ sequences. Folding in the HP model is NP-complete in general~\cite{berger1998protein}, and designability on lattice chains is known to be very uneven across structures~\cite{li1996emergence}, but at $N=16$ everything is enumerable and the success criterion---the ground state is unique and equals the target---is unambiguous.

We use this toy to compare, at an equal and explicit budget $K$ of oracle calls per design, a conditional autoregressive designer with simulated annealing and with two simple baselines, evaluated on target conformations never seen in training. Our questions are: (i) does the learned designer generalise to unseen conformations, (ii) does it beat search at equal oracle budget, (iii) is the conditioning on the target responsible, and (iv) how much do simple design choices of the baseline (objective, initialisation) change the answer? Our contributions are modest and specific: an exact-oracle protocol with a leakage-aware split, a five-seed comparison with ablations and a data-size sweep, and an honest report of where the learned model wins (small and medium budgets) and where it does not (a trivial heuristic at $K=1$; a better-posed annealing baseline at $K=1000$). All numbers come from a CPU-scale study of a single chain length and should not be read as a statement about real protein design.
